\documentclass[11pt]{article}

\usepackage[margin=0.6in]{geometry}
\usepackage[default]{sourcesanspro}
\usepackage[hidelinks]{hyperref}
\usepackage{enumitem}
\setlist[itemize]{left=0pt..1em, topsep=2pt, itemsep=2pt}
\usepackage{titlesec}

\titleformat{\section}{\large\bfseries\vspace{-4pt}}{}{0pt}{\vspace{-2pt}}
\titleformat{\subsection}{\bfseries\vspace{-3pt}}{}{0pt}{}

\begin{document}

\begin{flushright}
\today
\end{flushright}

\begin{center}
{\Large \textbf{Oth\'on Gonz\'alez}}\\[4pt]
\href{mailto:othon.gonzalezc@gmail.com}{othon.gonzalezc@gmail.com} \\
+52-55-3332-8699 \\
\end{center}

\section*{Professional Summary}
\vspace{-10pt}
\noindent\rule{\textwidth}{0.4pt}
Senior AI Engineer with experience across enterprise Agentic AI platforms, automotive Advanced Driver Assistance Systems (ADAS), and aerospace systems. Background in applied AI, LLM-based agents, computer vision, radar perception, model evaluation, observability, and safety-critical engineering environments. Most recently worked on production-oriented AI platforms involving agentic workflows, Retrieval-Augmented Generation (RAG), telemetry and tracing, data infrastructure, containerization, and cloud-native deployment. Experienced in building and validating AI systems across enterprise software, autonomous mobility, and aerospace applications.

\section{Skills}
\vspace{-10pt}
\noindent\rule{\textwidth}{0.4pt}
{\raggedright\setlength{\parindent}{0pt}
\textbf{Artificial Intelligence / Machine Learning:} Large Language Models (LLMs), Agentic AI workflows, Tool Calling, Agent evaluation, Retrieval-Augmented Generation (RAG), Computer Vision, PyTorch, TensorFlow, AI-assisted testing

\textbf{Data / Infrastructure:} PostgreSQL, Vector retrieval, Docker, Docker Compose, Kubernetes fundamentals, Cloud Computing, Observability fundamentals, Telemetry, Execution Tracing

\textbf{Programming:} Python, C++, MATLAB, Bash, Linux/GNU

\textbf{ADAS / Aerospace Systems:} Radar virtualization, Multi-sensor perception validation, Radar systems, Camera systems, LiDAR systems, dSPACE Vehicle Simulator, Safety-critical systems

\textbf{Software / Delivery:} Git, GitHub, JIRA, n8n, Agile, Scaled Agile Framework (SAFe)

\textbf{Engineering Tools:} SolidWorks, AutoCAD, Arduino, CNC, Laser cutting

\textbf{Languages:} Spanish (native), English (professional working proficiency)
\par}
\section{Experience}
\vspace{-10pt}
\noindent\rule{\textwidth}{0.4pt}

\subsection{Teradata}
\textbf{AI Engineer / Agentic AI Platform Engineer} \hfill Dec 2025 -- Jul 2026
\begin{itemize}
\item Contributing to an internal enterprise Agentic AI platform for deploying, consuming, evaluating, and observing AI agents.
\item Engineering LLM-based workflows with LangGraph, structured tool execution, reusable agent skills, MCP-based integrations, and Retrieval-Augmented Generation (RAG).
\item Designing the observability and distributed tracing layer using OpenTelemetry, capturing token usage, latency, cost, errors, model metadata, tool calls, execution paths, and agent outcomes.
\item Managing observability and monitoring components across Grafana, Prometheus, Loki, and distributed tracing pipelines to support debugging, performance analysis, and FinOps.
\item Building structured telemetry and evaluation pipelines using PostgreSQL/pgvector and DeepEval, enabling agent evaluation, analytics, dataset generation, and performance-driven model routing.
\item Developing containerized and cloud-native environments with Docker Compose and Kubernetes, supporting multi-service development, testing, and production-readiness.
\end{itemize}


\subsection{Continental Autonomous Mobility}
\textbf{Algorithms Test / Data Engineer (ADAS)} \hfill Jul 2022 -- Nov 2025
\begin{itemize}
\item Developed and implemented AI-driven testing tools for radar perception systems, reducing validation time for AI-based radar and camera model workflows.
\item Developed radar virtualization methods to support training and evaluation of AI models for perception-related use cases.
\item Engineered and deployed agentic AI and automated workflows using modern tooling, streamlining testing activities and improving team productivity.
\item Supported a strategic partnership with the National Polytechnic Institute through a mutual cooperation agreement, enabling collaborative research and development opportunities.
\item Designed and executed comprehensive test cases for multi-sensor Radar, Camera, and LiDAR perception pipelines to support system reliability.
\item Applied Agile and Scaled Agile Framework (SAFe) methodologies to coordinate complex projects within cross-functional teams.
\end{itemize}

\subsection{CentroGEO}
\textbf{Postdoctoral Project Assistant} \hfill Jul 2021 -- Jul 2023
\begin{itemize}
\item Researched and developed Natural Language Processing (NLP) models for generative tasks, including automated image captioning.
\item Worked with image-generative models, including DALL-E and Stable Diffusion, for research applications.
\item Managed data workflows for large-scale model training and inference, improving computational efficiency.
\item Architected and trained deep learning models, including Recurrent Neural Networks (RNNs), Long Short-Term Memory networks (LSTMs), and Transformers, for computer vision and language-related research problems.
\end{itemize}

\subsection{External Consultant}
\textbf{Safety Manager} \hfill Feb 2016 -- Jun 2018
\begin{itemize}
\item Led operational safety initiatives for a Mexican ground handling services provider.
\item Supported the successful certification of the Toluca station for International Standard for Business Aircraft Handling (IS-BAH) Safety Stages 1 and 2.
\item Rectified a critical permit renewal situation, negotiating with Mexican authorities to maintain essential aviation permits and prevent operational disruption.
\end{itemize}

\subsection{Aerom\'exico}
\textbf{Regional Airports Chief and Processes Analyst} \hfill Feb 2012 -- Oct 2013
\begin{itemize}
\item Served as liaison between regional airports and corporate headquarters, improving communication, process alignment, and operational follow-up.
\item Owned and improved the national Key Performance Indicator (KPI) for baggage irregularities, applying analytical methods to enhance tracking and reduce incidents.
\item Supported the successful renewal of International Air Transport Association Safety Audit for Ground Operations (ISAGO) certification and participated in multiple regulatory audits.
\end{itemize}

\subsection{ITESM / UVM / CentroGEO}
\textbf{Professor}
\begin{itemize}
\item Developed and taught undergraduate and graduate curricula in engineering subjects including Statistics, Aerodynamics, Physics, and specialized aeronautical topics.
\end{itemize}

\section{Education}
\vspace{-10pt}
\noindent\rule{\textwidth}{0.4pt}
\textbf{PhD in Computer Vision} | CICATA Quer\'etaro, IPN \hfill Jan 2020 \\
Thesis: Precise remote measurement of thermal radiance dynamics using UAVs \\
Final Grade: 9.0/10 

\vspace{8pt}
\noindent
\textbf{MSc in Aerodynamic Optimization} | CICATA Quer\'etaro, IPN \hfill Feb 2016 \\
Thesis: Data-driven morphable wing for tracking and mapping using small UAVs \\
Final Grade: 9.5/10 \hfill \textit{Cum Laude}

\vspace{8pt}
\noindent
\textbf{Eng in Aeronautics} | ESIME Ticom\'an, IPN \hfill Mar 2010 \\
Thesis: Software tool to determine the capacity of different runway systems \\
Final Grade: 9.4/10 \hfill \textit{Cum Laude}

\section{Research \& Publications}
\vspace{-10pt}
\noindent\rule{\textwidth}{0.4pt}
\begin{itemize}
\item 2019. Gonz\'alez-Ch\'avez, O., et al. Radiometric Calibration of Infrared Thermal Cameras. IEEE Transactions on Instrumentation and Measurement.
\item 2019. Gonz\'alez-Ch\'avez, O., et al. Thermal Radiation Dynamics of Soil Surfaces with UAVs. Mexican Conference on Pattern Recognition.
\item 2023. Moctezuma, D., et al. Video Captioning: a comparative review. Computer Vision and Image Understanding, Vol. 231.
\item 2024. Gonz\'alez-Ch\'avez, O., et al. Are metrics measuring what they should? Evaluation of image captioning task metrics. Signal Processing: Image Communication, Vol. 120.
\end{itemize}

\end{document}